%==============================================================================
% MATH 5080
% Class:        8
% Date:         Sept 17th
%==============================================================================
\begin{proof}[Proof of 
    \ref{theorem:bdd-implies-cvgt-subseq} 
    $\implies$
    \ref{theroem:bdd-dvgt-implies-multiple-cvgt-subseq}]
    \hfill

    Fix a bounded and divergent sequence $\{ a_n \}_n $
    \ref{theorem:bdd-implies-cvgt-subseq},
    $\exists \{ a_{N_k} \}_k $ and $L_1 \in \mathbb{R}$  s.t. 
    $a_{N_K} \to L$.
    By definition of divergence, we know that $ \{ a_n \} $ 
    does not converge, i.e. fix $\ep > 0$,
    \begin{align*}
        & \forall N > 0, \exists n > N \text{ s.t. } 
        |a_n - L_1| > \ep \\
        \implies & \exists \{ a_{n_l} \}_l \text{ s.t. }
        |a_{n_l} - L_1| > \ep, \forall l \in \mathbb{N}
    \end{align*}
    Since $\{ a_{n_l} \}_l $ bounded, it has a convergent subsequence by 
    \ref{theorem:bdd-implies-cvgt-subseq}. WLOG, assume $\{ a_{n_l} \} $
    converges to $L_2$. 
    \medskip

    Since $|a_{n_l} - L_2| > \ep, \forall l \in \mathbb{N}$, 
    $|L_2 - L_1| \geq \ep > 0$.
\end{proof}

\begin{proof}[Proof of \ref{theorem:bdd-implies-cvgt-subseq}]
    \hfill

    Consider 2 cases 
    \begin{itemize}
        \item Case 1: $\{ a_n \} $ finite
        Then at least one value is attained infinitely many times, 
        take a subsequence consisting only of that value, and
        we have found a convergent subsequence. 
        \item Case 2: $\{ a_n \} $ infinite
        By BW, the set of values attained by $\{ a_n \} $
        has an accumulation point say $L$.
        \medskip

        We want to construct $\{ a_[n_k] \} $ converging to $L$. 
        \begin{itemize}
            \item Let $\ep_1 \coloneq 1$, and 
            $a_{n_1} \in N^*(L, \ep_1) \cap \{ a_n \}$.
            \item Let $\ep_2 = \frac{1}{2}$. 
            Since $N^*(L, \ep_2) \cap \{ a_n \}_{n \geq n_1 + 1} $ has infinitely
            many points, $\exists n_2 > n$ s.t. $a_{n_2} \in N^*(L, \ep_2)$. 
            \item Repeat this process, if $a_{n_1}, a_{n_2}, \hdots a_{n_k}$ chosen, 
            let $\ep_{k + 1} = \frac{1}{k + 1}$ and pick
            \[
                a_{n_{k + 1}} \in N^*(L, \ep_{k + 1}) \cap \{ a_n \}_{n \geq n_K + 1}
            \], 
            and we have constructed a subsequence $\{ a_{n_k} \} $ such that 
            $a_{n_k} \rightarrow L$
        \end{itemize} 
    \end{itemize} 
\end{proof}

\begin{definition}[Subsequential limits]
    Let $\{ a_n \} \subset \mathbb{R}$, the set of subsequential limits of 
    $\{ a_n \} $ is the set of $x \in \mathbb{R} \cup \{ \pm \infty \} $ such that 
    some subsequence $\{ a_{n_k} \} $ converges to $x$.
\end{definition}

\begin{remark}
    $\{ a_n \} $ bounded if and only if its set of subsequential limits 
    is a nonempty and bounded subset of $\mathbb{R}$.
\end{remark}

\begin{definition}[Extended reals]
    In the extended real number system $R \cup \{ \pm \infty \} $,
    \begin{itemize}
        \item we have order
            \[
                -\infty < x < \infty, \forall x
            \]
        \item we do not have a field 
        \item we have the following conventions, $\forall x \in \mathbb{R}$
        \begin{align*}
            & x + \infty = \infty \\
            & x - \infty =  - \infty \\
            & \frac{x}{+ \infty} = \frac{x}{- \infty} = 0 \\
            &x > 0 \implies x \times (\pm \infty) = \pm \infty \\
            &x < 0 \implies x \times (\pm \infty) = \mp \infty \\
            \end{align*}
    \end{itemize} 
\end{definition}


\begin{definition}[Limit superior, limit inferior]
    The limit superior of $\{ a_n \} $ is defined as the supremum 
    of the set of subsequential limit, denoted as 
    \[
        \limsup_{ n \to \infty} a_n \in \mathbb{R} \cup \{ \pm \infty \} 
    \]

    Similarly, the limit inferior is defined as the 
    infimum of the set of subsequential limits, denoted
    $\liminf_{n \to \infty} a_n$.
\end{definition}

\begin{remark}
    The limit inferiors and superiors always exists for any sequence 
    in $\mathbb{R} \cup \{ \pm \infty \} $.
    \begin{itemize}
        \item $\{ a_n \} $ bounded from above $\iff$ 
        $\limsup_{n \to \infty} a_n \in \mathbb{R} \cup \{ -\infty \} $. 
        Otherwise $\limsup_{n \to \infty} a_n = \infty$. 
        \item $\{ a_n \} $ bounded from bleow $\iff$ 
        $\liminf_{n \to \infty} a_n \in \mathbb{R} \cup \{ \infty \} $. 
        Otherwise $\liminf{n \to \infty} a_n = -\infty$. 
        \item $\{ a_n \} $ bounded $\implies$
        $\liminf a_n, \limsup a_n \in \mathbb{R}$.
    \end{itemize} 
\end{remark}


\begin{proposition}
    \hfill
    \begin{enumerate}
        \item $\liminf_{n \to \infty} a_n \leq \limsup_{n \to \infty}a_n$
        \item A sequence converges if and only if its limsup and liminf agree
        \[
            \lim_{n \to \infty} a_n = L \in \mathbb{R} \cup \{ \pm \infty \} 
            \iff 
            \liminf_{n \to \infty} a_n = 
            \limsup_{n \to \infty} a_n = 
            L
        \]
    \end{enumerate}
\end{proposition}
\begin{proof}
    \textbf{
        WTS
        $\lim_{n \to \infty} a_n = L$
        $\implies$
        $\liminf_{n \to \infty} a_n = \limsup_{n \to \infty} a_n = L$
    }
    Since every subsequence of a convergent sequence converges to the 
    same limit.
    \medskip

    \textbf{
        WTS
        $\liminf_{n \to \infty} a_n = \limsup_{n \to \infty} a_n = L$
        $\implies$
        $\lim_{n \to \infty} a_n = L$
    }
    ATC otherwise, i.e. 
    \begin{align*}
        & \lim_{n \to \infty} a_{n} \neq L 
        \implies & \exists \{ a_{n_k} \} \text{ s.t. }
        \exists \ep_0, |a_{n_k} - L | > \ep_0, \forall k \in \mathbb{N}
    \end{align*}
    Since $\liminf a_n = \limsup a_n = L$, $\{ a_n \} $ 
    must be bounded. Hence $ \{ a_{n_k} \} $ also bounded.
    Hence $\{ a_{n_k} \} $ has a convergent subsequence,
    which we still denote $\{ a_{n_k} \} $. 
    Since this is also a subsequence of the original $\{ a_n \} $, we know 
    that $\lim_{k \to \infty}a_{n_k} = L$. However, this leads
    to a contradiction because take $\ep = \frac{\ep_0}{2}$, then
    $\forall k, |a_{n_k} - L | > \ep_0 > \ep$.
\end{proof}


\begin{proposition}
    If $\{ a_n \} $ is bounded then its set of subsequential limits is closed. 
\end{proposition}

\begin{proof}
    \textit{
        The idea is to take any accumulation point
        in the set of subsequential limits of $a_n$,
        and show that some subsequence indeed
        converges to that accumulation point.
    }
    \medskip

    If $\{ a_n \} $ bounded, 
    then all of its subsequential limits are in $\mathbb{R}$.
    Let $E$ denote the set of subsequential limits. 
    WTS that $E$ contains all of its accumulation points, 
    i.e. $\forall p \in E^{\prime}, p \in E$. 
    \medskip

    Now we want to construct a subsequence that 
    converges to $p$
    \begin{itemize}
        \item First choose $p_1 \in N^*(p, 1) \cap E$. 
    Let $\ep_1 = 1 - |p-p_1| > 0$.
    $\exists a_{n_1}$ s.t. $|a_{n_1} - p_1| < \ep_1 \implies$
    $|a_{n_1} - p| < 1$.
        \item $\vdots$
        \item 
    Suppose $a_{n_1} , \hdots a_{n_{k - 1}}$ are found, s.t. 
    $n_1 < n_2 < \hdots n_{k-1}$,
    choose $p_k \in N^*(p, \frac{1}{k}) \cap E$. 
    Let $\ep_k = \frac{1}{k} - |p - p_k| > 0$.
    $\exists a_{n_k}$ s.t. $n_k > n_{k-1}$ and 
    $|a_{n_k} - p_k| < \ep_k$. 
    Hence $|a_{k_k} - p| < \frac{1}{k}$. 
    \end{itemize} 

    Therefore $a_{n_k} \to p$ and $p \in E$.
\end{proof}



\begin{theorem}
    \label{theorem:limsup-properties}
    If $\{ a_n \} $ is bounded and $x \in \mathbb{R}$, TFAE 
    \begin{enumerate}
        \item $\limsup_{n \to \infty} a_n = x$ 
        \item 
        \begin{enumerate}
            \item $\forall \ep > 0, \exists N > 0$ s.t. 
            $\forall n > N$, $a_n < x + \ep$ AND
            \item $\forall \ep > 0, \forall N > 0$,
            $\exists n > N$, s.t. $a_n > x - \ep$
        \end{enumerate}
        \item \begin{enumerate}
            \item There is $\{ a_{n_k} \} $ that converges to $x$ 
            AND
            \item If $\{ a_{n_j} \} $ converges, then 
            $\lim_{j \to \infty} a_j \leq x$
        \end{enumerate}
        \item Let $b_n = \sup_{m \geq n} a_m$, then $b_n$ converges
        and $\lim_{n \to \infty} b_n = x$
    \end{enumerate}
\end{theorem}

% \begin{proof}[Proof of 
%         \ref{theorem:limsup-properties},
%         1 $\iff$ 2
%     ]
%     \begin{itemize}
%         \item 1 $\implies$ 2(a). 
%         \medskip

%         ATC otherwise, fix $\ep > 0$, suppose 
%         $\not \exists N > 0$ s.t. $\forall n > N$, $a_n < x + \ep$.  Then $\exists \{ a_{n_k} \}_k $ s.t. $a_{n_k} \geq x + \ep, \forall k$. 
%         By \ref{theorem:bdd-implies-cvgt-subseq}, it has a subsequence 
%         that converges to some $L \geq x + \ep$, hence a contradiction. 
%         \item 1 $\implies$ 2(b)
%         \medskip

%         ATC otherwise, then $\exists N > 0$ s.t. $\forall n > N$, 
%         $a_n \leq x - \ep$. Hence all convergent subsequences have limits $x < \ep$< 
%         \medskip 

%         Hence $x \leq x - \ep$, which leads to a contradiction.
%         \item 2 $\implies$ 1.
%         \medskip

%         Fix $\ep > 0$, by \ref{theorem:limsup-properties},
%         2.(a), all subsequential limits are upper bound by 
%         $x + \ep$.
%         \[
%             \limsup a_n \leq x + \ep \implies \limsup a_n \leq x + \ep
%         \]
%         We claim that $\limsup a_n = x$. If not, then $\limsup (a_n) < x$< 
%         le $\ep = \frac{1}{2}(x - \limsup(a_n)) > 0$. By 
%         2.(b), we can find a subsequence $\{ a_{n_k} \} $ s.t. 
%         $a_{n_k} > x - \ep$. By
%         \ref{theroem:bdd-dvgt-implies-multiple-cvgt-subseq},
%         WLOG, assume $a_{n_k} \to L$, then 
%         \[
%             L \geq x - \ep > \limsup a_N
%         \]
%         which is a contradiction.
%     \end{itemize}
% \end{proof}

% \remind{fill in additional material from class 8}

\begin{proof}[Proof of 
        \ref{theorem:limsup-properties},
    ]
    In general, let $E$ denote the 
    set of subsequential limits of $a_n$.
    \begin{itemize}
        \item (1) $\iff$ (2)
        \begin{enumerate}
            \item (1) $\implies$ (2)
            \begin{itemize}
                \item (1) $\implies$ (2a): ATC otherwise,
                then $\exists$ infinitely many terms
                where $a_n \geq x + \ep$.
                Then they form a subsequence $ \{ a_{n_k} \} $
                which is also bounded. 
                Hence this subsequence itself has a convergent
                subsequence converging to some
                limit $L \geq x + \ep > x$. 
                Hence contradicting the assumption
                that $x$ is a supremum.
                \item (1) $\implies$ (2b)
                ATC otherwise, then only finitely many terms
                $a_n > x - \ep$. Hence,
                $\exists N >0$ s.t. $\forall n > N$,
                $a_n \leq x - \ep$. Hence 
                all subsequential limits must be less 
                than $x - \ep$. This contradicts
                the assumption that $x$ is a minimal upper bound.
            \end{itemize} 
            \item (1) $\impliedby$ (2): Fix any $\ep > 0$, 
            $\exists N \text{ s.t. } \forall n > N$,
            $a_n < x + \ep$.
            Hence any subsequential limit is smaller than $x + \ep$.
            Therefore $\limsup a_n \leq x + \ep$. Since
            $\ep$ is arbitrary, $\limsup a_n \leq x$.
        \end{enumerate}
        \item (1) $\iff$ (3)
        \begin{enumerate}
            \item (1) $\implies$ (3)
            Since the set of subsequential limits is closed,
            the supremum is in $E$. Hence
            some subsequence converges to it. 
            By definition of supremum, $x \geq p, \forall p \in E$.
            \item (1) $\impliedby$ (3): by definition
        \end{enumerate} 
        \item (1) $\iff$ (4)
        \begin{enumerate}
            \item (1) $\implies$ (4): 
            Since $\{ a_n \} $ bounded, 
            $\{ \sup_{m \geq n} a_m\}_n $ is also bounded
            by the same bounds. 
            Also note
            \begin{align*}
                \sup_{m \geq n + 1} a_m \leq
                \sup_{m \geq n} a_m
            \end{align*}
            Hence the sequence is decreasing. 
            By MCT
            \begin{align*}
                \lim_{n \to \infty}
                \sup_{m \geq n} a_m = L \in \mathbb{R}
            \end{align*}
            It suffices to show that $L = x$. 
            Fix $\ep > 0$. 
            \begin{itemize}
                \item By (2a): $\exists N > 0$ s.t. $\forall n > N$
                $a_n < x + \frac{\ep}{2}$. Hence 
                \[
                    b_n = \sup_{m \geq n} a_m \leq 
                    x + \frac{\ep}{2} 
                    < x + \ep
                \]
                Hence $L = \lim_{n \to \infty}b_n \leq x + \ep$
                $\implies$
                $L \leq x$
                \item By (2b): $\forall n > 0$, 
                \begin{align*}
                    b_n = \sup_{m \geq n} a_m > x  - \ep
                \end{align*}
                Hence $L = \lim_{n \to  \infty}b_n \geq x - \ep$
                $\implies$
                $L \geq x$.
            \end{itemize} 
            \item (1) $\impliedby$ (4):
            Prove by showing (4) $\implies$ (3).
            \medskip

            We pick 
            \begin{itemize}
                \item $n_1$ s.t. $a_{n_1} \geq \sup_{m \geq 1} a_m - 1$
                \item $n_2$ s.t. $n_2 > n_1$ and 
                $a_{n_2} \geq \sup_{m \geq n_1 + 1} a_m - \frac{1}{2}$.
                \item $\vdots$
                \item $n_k$ s.t. $n_k > n_{k-1}$ s.t. $a_{n_k} \geq$
                $\sup_{m \geq n_{k-1} + 1} a_m - \frac{1}{k}$.
            \end{itemize} 
            Let $L = \lim_{n \to \infty} \sup_{m \geq n} a_m$.
            \medskip

            WTS $\lim_{n \to \infty}a_{n_k} = L$.
            Fix $ \ep > 0$. Choose $N > 0$ s.t. 
            $\forall n > N$,
            \[
                \left| \sup_{m \geq n} a_m - L\right|  < \frac{\ep}{2}
            \]
            Then take $N_0 = \max \{ N + 1, \frac{2}{\ep} \} $,
            one has $\forall k > N_0$, 
            \begin{align*}
                \left|a_{n_k}- L \right| &\leq 
                \left| a_{n_k} - \sup_{m \geq n_{k -1} + 1} a_m \right|  
                +
                \left| \sup_{m \geq n_{k -1} + 1} - L\right|  \\
                & \leq \frac{1}{k} + \frac{\ep}{2} \\
                & \leq \frac{\ep}{2} + \frac{\ep}{2} = \ep
            \end{align*}
            Hence we have shown (3a). 
            \medskip

            Let $\{ a_{n_k} \} $ be any convergent subsequence.
            of $ \{ a_n \} $. We have
            \begin{align*}
                &a_{n_k} \leq \sup_{m \geq n_k} a_m, \forall k \\
                \implies & 
                \lim_{k \to \infty} a_{n_k} \leq 
                \lim_{k \to \infty} \sup_{m \geq n_k} a_m = L
            \end{align*}
        \end{enumerate}
    \end{itemize}
\end{proof}





